% ============================================================
\chapter*{参考文献}
\addcontentsline{toc}{chapter}{参考文献}

{\small
\begin{itemize}[leftmargin=1.2em,itemsep=0.25em,label={}]
\item Allaj, E. (2013). The Black--Litterman model: a consistent estimation of the parameter tau. \textit{Financial Markets and Portfolio Management} 27(2), 217--251.
\item Almgren, R. and N. Chriss (2006). Optimal portfolios from ordering information. \textit{Journal of Risk} 9(1), 1--47.
\item Alpern, A. and S. T. Rachev (2025). Black Litterman and ESG portfolio optimization. arXiv:2511.21850.
\item Artzner, P., F. Delbaen, J.-M. Eber and D. Heath (1999). Coherent measures of risk. \textit{Mathematical Finance} 9(3), 203--228.
\item Asness, C. S., T. J. Moskowitz and L. H. Pedersen (2013). Value and momentum everywhere. \textit{Journal of Finance} 68(3), 929--985.
\item Baker, S. R., N. Bloom and S. J. Davis (2016). Measuring economic policy uncertainty. \textit{Quarterly Journal of Economics} 131(4), 1593--1636.
\item Barroso, P. and P. Santa-Clara (2015). Momentum has its moments. \textit{Journal of Financial Economics} 116(1), 111--120.
\item Beach, S. L. and A. G. Orlov (2007). An application of the Black--Litterman model with EGARCH-M-derived views for international portfolio management. \textit{Financial Markets and Portfolio Management} 21(2), 147--166.
\item Bessler, W., H. Opfer and D. Wolff (2017). Multi-asset portfolio optimization and out-of-sample performance: an evaluation of Black--Litterman, mean-variance, and na\"ive diversification approaches. \textit{European Journal of Finance} 23(1), 1--30.
\item Bevan, A. and K. Winkelmann (1998). Using the Black-Litterman global asset allocation model: three years of practical experience. Goldman Sachs, Fixed Income Research.
\item Black, F. and R. Litterman (1990). Asset allocation: combining investor views with market equilibrium. Goldman Sachs, Fixed Income Research.
\item Black, F. and R. Litterman (1992). Global portfolio optimization. \textit{Financial Analysts Journal} 48(5), 28--43.
\item Blamont, D. and N. Firoozye (2003). Asset allocation model. Deutsche Bank, Global Markets Research.
\item Brogaard, J. and A. Detzel (2015). The asset-pricing implications of government economic policy uncertainty. \textit{Management Science} 61(1), 3--18.
\item Caldara, D. and M. Iacoviello (2022). Measuring geopolitical risk. \textit{American Economic Review} 112(4), 1194--1225.
\item Carhart, M. M. (1997). On persistence in mutual fund performance. \textit{Journal of Finance} 52(1), 57--82.
\item Chen, N.-F., R. Roll and S. A. Ross (1986). Economic forces and the stock market. \textit{Journal of Business} 59(3), 383--403.
\item Cheung, W. (2013). The augmented Black--Litterman model: a ranking-free approach to factor-based portfolio construction and beyond. \textit{Quantitative Finance} 13(2), 301--316.
\item Colasanto, F., L. Grilli, D. Santoro and G. Villani (2022). BERT's sentiment score for portfolio optimization: a fine-tuned view in Black and Litterman model. \textit{Neural Computing and Applications} 34(20), 17507--17521.
\item Daniel, K. and T. J. Moskowitz (2016). Momentum crashes. \textit{Journal of Financial Economics} 122(2), 221--247.
\item Duqi, A., L. Franci and G. Torluccio (2014). The Black--Litterman model: the definition of views based on volatility forecasts. \textit{Applied Financial Economics} 24(19), 1285--1296.
\item Engle, R. F., D. M. Lilien and R. P. Robins (1987). Estimating time varying risk premia in the term structure: the ARCH-M model. \textit{Econometrica} 55(2), 391--407.
\item Fabozzi, F. J., S. M. Focardi and P. N. Kolm (2006). Incorporating trading strategies in the Black--Litterman framework. \textit{Journal of Trading} 1(2), 28--37.
\item Fama, E. F. and K. R. French (1993). Common risk factors in the returns on stocks and bonds. \textit{Journal of Financial Economics} 33(1), 3--56.
\item Fama, E. F. and K. R. French (2015). A five-factor asset pricing model. \textit{Journal of Financial Economics} 116(1), 1--22.
\item Florencio, M. (2026). Grounding investor views: neural predicates in the Black-Litterman model. arXiv:2607.20533.
\item Giacometti, R., M. Bertocchi, S. T. Rachev and F. J. Fabozzi (2007). Stable distributions in the Black--Litterman approach to asset allocation. \textit{Quantitative Finance} 7(4), 423--433.
\item Grinold, R. C. (1994). Alpha is volatility times IC times score. \textit{Journal of Portfolio Management} 20(4), 9--16.
\item Grinold, R. C. and R. N. Kahn (1999). \textit{Active Portfolio Management}, 2nd ed. McGraw-Hill.
\item Han, Y. and J. Li (2023). The impact of global economic policy uncertainty on portfolio optimization: A Black--Litterman approach. \textit{International Review of Financial Analysis} 86, 102476.
\item Hassan, T. A., S. Hollander, L. van Lent and A. Tahoun (2019). Firm-level political risk: measurement and effects. \textit{Quarterly Journal of Economics} 134(4), 2135--2202.（この資料では NBER の版を読んだ）
\item He, G. and R. Litterman (1999). The intuition behind Black-Litterman model portfolios. Goldman Sachs, Investment Management Research.
\item Herold, U. (2003). Portfolio construction with qualitative forecasts. \textit{Journal of Portfolio Management} 30(1), 61--72.
\item Idzorek, T. M. (2007). A step-by-step guide to the Black-Litterman model: incorporating user-specified confidence levels. In S. Satchell (ed.), \textit{Forecasting Expected Returns in the Financial Markets}, Academic Press, 17--38.（この資料では 2004 年 7 月 20 日の版を読んだ）
\item Jegadeesh, N. (1990). Evidence of predictable behavior of security returns. \textit{Journal of Finance} 45(3), 881--898.
\item Jegadeesh, N. and S. Titman (1993). Returns to buying winners and selling losers: implications for stock market efficiency. \textit{Journal of Finance} 48(1), 65--91.
\item Jorion, P. (2003). Portfolio optimization with tracking-error constraints. \textit{Financial Analysts Journal} 59(5), 70--82.
\item Kadavath, S. ほか (2022). Language models (mostly) know what they know. arXiv:2207.05221.
\item Kolm, P. N. and G. Ritter (2017). On the Bayesian interpretation of Black--Litterman. \textit{European Journal of Operational Research} 258(2), 564--572.
\item Kolm, P. N., G. Ritter and J. Simonian (2021). Black--Litterman and beyond: the Bayesian paradigm in investment management. \textit{Journal of Portfolio Management} 47(5).
\item Ledoit, O. and M. Wolf (2004). A well-conditioned estimator for large-dimensional covariance matrices. \textit{Journal of Multivariate Analysis} 88(2), 365--411.
\item Lee, W. (2000). \textit{Advanced Theory and Methodology of Tactical Asset Allocation}. Wiley.
\item Lee, Y., Y. Kim, J. Kim, S. Kim and Y. Lee (2025). LLM-enhanced Black-Litterman portfolio optimization. CIKM 2025 Workshop on FinAI. arXiv:2504.14345（第 2 版）.
\item Lehmann, B. N. (1990). Fads, martingales, and market efficiency. \textit{Quarterly Journal of Economics} 105(1), 1--28.
\item Litterman, R. (2003). \textit{Modern Investment Management: An Equilibrium Approach}. Wiley.
\item Mantshimuli, L. A. and J. W. Muteba Mwamba (2025). Enhancing portfolio optimization with multi-LLM sentiment aggregation: A Black-Litterman integration approach. \textit{Investment Management and Financial Innovations} 22(3), 213--226.
\item Meucci, A. (2005). \textit{Risk and Asset Allocation}. Springer.
\item Meucci, A. (2006). Beyond Black-Litterman: views on non-normal markets. \textit{Risk} 19(2), 87--92.
\item Meucci, A. (2006). Beyond Black-Litterman in practice: a five-step recipe to input views on non-normal markets. \textit{Risk} 19(9), 114--119.
\item Meucci, A. (2008). Fully flexible views: theory and practice. \textit{Risk} 21(10), 97--102.
\item Meucci, A. (2009). Enhancing the Black--Litterman and related approaches: views and stress-test on risk factors. \textit{Journal of Asset Management} 10(2), 89--96.
\item Meucci, A. (2010). The Black-Litterman approach: original model and extensions. \textit{The Encyclopedia of Quantitative Finance}, Wiley. SSRN 1117574（この資料では 2010 年 10 月 12 日の版を読んだ）.
\item Michaud, R. O. (1989). The Markowitz optimization enigma: is `optimized' optimal? \textit{Financial Analysts Journal} 45(1), 31--42.
\item Nelson, D. B. (1991). Conditional heteroskedasticity in asset returns: a new approach. \textit{Econometrica} 59(2), 347--370.
\item O'Toole, R. (2017). The Black--Litterman model: active risk targeting and the parameter tau. \textit{Journal of Asset Management} 18(7).
\item P\'astor, \v{L}. and P. Veronesi (2012). Uncertainty about government policy and stock prices. \textit{Journal of Finance} 67(4), 1219--1264.
\item P\'astor, \v{L}. and P. Veronesi (2013). Political uncertainty and risk premia. \textit{Journal of Financial Economics} 110(3), 520--545.
\item Pedersen, L. H., S. Fitzgibbons and L. Pomorski (2021). Responsible investing: the ESG-efficient frontier. \textit{Journal of Financial Economics} 142(2), 572--597.
\item Qian, E. and S. Gorman (2001). Conditional distribution in portfolio theory. \textit{Financial Analysts Journal} 57(2), 44--51.
\item Rockafellar, R. T. and S. Uryasev (2000). Optimization of conditional value-at-risk. \textit{Journal of Risk} 2(3), 21--41.
\item Roll, R. (1992). A mean/variance analysis of tracking error. \textit{Journal of Portfolio Management} 18(4), 13--22.
\item Satchell, S. and A. Scowcroft (2000). A demystification of the Black--Litterman model: managing quantitative and traditional portfolio construction. \textit{Journal of Asset Management} 1(2), 138--150.
\item Sharpe, W. F. (1974). Imputing expected security returns from portfolio composition. \textit{Journal of Financial and Quantitative Analysis} 9(3), 463--472.
\item Spears, T., S. Zohren and S. Roberts (2023). View fusion vis-\`a-vis a Bayesian interpretation of Black-Litterman for portfolio allocation. arXiv:2301.13594.
\item Sun, L., B. Fu and W. Liu (2026). Investor sentiment and stock portfolio construction: A study based on the Black-Litterman model. \textit{PLOS ONE} 21(9).
\item Sun, Y., Y. Wu and G. De (2023). A novel Black-Litterman model with time-varying covariance for optimal asset allocation of pension funds. \textit{Mathematics} 11(6), 1476.
\item Teplova, T., E. Mikova, Q. Munir and N. Pivnitskaya (2023). Black-Litterman model with copula-based views in mean-CVaR portfolio optimization framework with weight constraints. \textit{Economic Change and Restructuring} 56(1), 515--535.
\item Theil, H. and A. S. Goldberger (1961). On pure and mixed statistical estimation in economics. \textit{International Economic Review} 2(1), 65--78.
\item Walters, J. (2014). The Black-Litterman model in detail. SSRN 1314585.
\item Xiong, M. ほか (2024). Can LLMs express their uncertainty? An empirical evaluation of confidence elicitation in LLMs. ICLR 2024.
\item Zuber (2017). 修士論文（ESG の点数で共分散の構造を変える BL）. University of Zurich, Center for Sustainable Finance and Private Wealth.
\item \textit{Incorporating Black-Litterman views in portfolio construction when stock returns are a mixture of normals} (2017). arXiv:1706.03287.
\end{itemize}
}

\section*{データ}

{\small
\begin{itemize}[leftmargin=1.2em,itemsep=0.25em,label={}]
\item Kenneth R. French, Data Library（Fama/French 5 Factors (2x3)、Momentum Factor（月次・日次））。
\item Economic Policy Uncertainty Index（Baker, Bloom and Davis）、policyuncertainty.com。
\item Geopolitical Risk Index（Caldara and Iacoviello）、matteoiacoviello.com。
\item 本プロジェクトの DB（\code{data/blm\_2015.sqlite}）。株価は yfinance、構成銘柄は Wikipedia、財務諸表と発行済株式数は SEC EDGAR から作った（『総合報告書』）。
\end{itemize}
}
